\documentclass[11pt,a4paper]{article}

\usepackage[a4paper,margin=1in]{geometry}
\usepackage{fontspec}
\IfFontExistsTF{Times New Roman}
  {\setmainfont{Times New Roman}}
  {\IfFontExistsTF{TeX Gyre Termes}{\setmainfont{TeX Gyre Termes}}{\setmainfont{Times}}}
\usepackage{newtxmath}
\usepackage{amsmath}
\usepackage{booktabs}
\usepackage{longtable}
\usepackage{array}
\usepackage{enumitem}
\usepackage{xcolor}
\usepackage{graphicx}
\usepackage{float}
\usepackage{listings}
\usepackage[round,authoryear]{natbib}
\usepackage[hidelinks]{hyperref}
\usepackage{xurl}
\usepackage{fancyhdr}
\usepackage{microtype}

\setlength{\parindent}{0pt}
\setlength{\parskip}{0.65em}
\setlist[itemize]{leftmargin=1.4em,itemsep=0.2em,topsep=0.2em}
\setlist[enumerate]{leftmargin=1.6em,itemsep=0.2em,topsep=0.2em}

\pagestyle{fancy}
\fancyhf{}
\lhead{\small Ticket 05 Study Report}
\rhead{\small American Put Validation}
\cfoot{\thepage}
\renewcommand{\headrulewidth}{0.3pt}

\lstset{
  basicstyle=\ttfamily\small,
  breaklines=true,
  frame=single,
  columns=fullflexible,
  keepspaces=true,
  rulecolor=\color{black!30},
  xleftmargin=0.5em,
  xrightmargin=0.5em
}

\newcommand{\code}[1]{\texttt{#1}}
\newcommand{\file}[1]{\path{#1}}
\newcolumntype{L}[1]{>{\raggedright\arraybackslash}p{#1}}
\newcolumntype{P}[1]{>{\raggedright\arraybackslash\ttfamily\footnotesize}p{#1}}

\title{\textbf{Ticket 05 Study Report: American Put Validation}\\
\large SURF2026 American Option Risk Surfaces}
\author{Codex-assisted implementation report}
\date{June 16, 2026}

\begin{document}
\maketitle

\begin{abstract}
Ticket 05 uses the Ticket 04 American CN/PSOR solver to validate American put behavior in a
representative Black-Scholes-with-dividends case. The ticket produces validation tables, selected
spot checks, obstacle and convergence summaries, and a value-versus-payoff figure. This is a
validation ticket, not a new solver ticket. It does not implement boundary extraction, Greeks,
stress maps, datasets, neural-network files, or American call validation.
\end{abstract}

\tableofcontents
\newpage

\section{Role of Ticket 05 in the Whole Solver Project}

Tickets 01 through 04 built the solver foundation in layers. Ticket 01 implemented payoff and
European Black-Scholes utilities. Ticket 02 implemented finite-difference grids and the
Black-Scholes spatial operator. Ticket 03 validated European Crank-Nicolson time stepping against
closed-form European prices. Ticket 04 added the PSOR/LCP core needed for the American payoff
obstacle.

Ticket 05 is the first validation ticket for a genuinely American option case. It does not add a
new pricing algorithm. Instead, it uses the Ticket 04 American CN/PSOR solver and asks whether the
American put outputs look financially and numerically credible in a representative setup.

The role of this ticket is evidence gathering:
\begin{itemize}
  \item check that computed American put values remain above payoff;
  \item check that computed American put values are at least European put values;
  \item summarize PSOR convergence metadata;
  \item inspect selected spot prices in the main region of interest;
  \item compare medium and fine grids lightly, without claiming full grid sensitivity.
\end{itemize}

This ticket intentionally stops before free-boundary extraction. It may report the obstacle gap
\(U-\Phi\), because that is a validation diagnostic. It does not turn that gap into a boundary
location.

\section{Theory and Methodology Connection}

The documented source basis for Ticket 05 is internal:
\begin{itemize}
  \item \file{reports/01_literature/literature_map.md};
  \item \file{reports/02_math/formulation_note.md};
  \item \file{reports/03_solver/solver_validation_plan.md};
  \item \file{reports/03_solver/tickets/ticket_04_psor_lcp_core.tex}.
\end{itemize}

The literature map identifies American option pricing as an obstacle problem. The formulation note
expresses the American option conditions as a variational inequality and then as a linear
complementarity problem after discretization. The solver validation plan states that American put
validation should include obstacle checks, continuation-premium or obstacle-gap sanity checks,
selected prices, and convergence evidence.

This reference integration pass adds verified citations for the external methodology behind the
internal source basis. American put finite-difference validation is connected to Brennan-Schwartz
\citep{brennan_schwartz_1977}; the LCP and obstacle language is connected to standard LCP
methodology \citep{cottle_pang_stone_2009}; and the European comparison benchmark rests on the
Black-Scholes/Merton framework \citep{black_scholes_1973,merton_1973}.

\section{Why American Put Validation Follows the PSOR/LCP Core}

A European validation run checks the continuation PDE. An American put validation run checks the
continuation PDE plus the early-exercise constraint. The American constraint is
\begin{equation}
  U(S,\tau) \geq \Phi_{\mathrm{put}}(S),
  \qquad
  \Phi_{\mathrm{put}}(S)=\max(K-S,0).
\end{equation}

Ticket 04 built the PSOR machinery that enforces this constraint at each Crank-Nicolson step.
Ticket 05 comes next because the solver core must now be used in a realistic put case, with
reported evidence rather than only smoke tests.

For beginner researchers, the distinction is important:
\begin{itemize}
  \item Ticket 04 asks: can the code enforce the obstacle in a reusable solver core?
  \item Ticket 05 asks: when that core is used for an American put, do the prices and diagnostics
        obey the basic financial checks?
\end{itemize}

\section{American Put Financial Intuition}

An American put gives the holder the right to sell the underlying asset for strike \(K\) at any
time up to maturity. When the spot \(S\) is low, the payoff \(K-S\) is high. Early exercise can
therefore be attractive for puts, especially when the put is deep in the money.

The American put should be at least as valuable as:
\begin{itemize}
  \item immediate exercise payoff, because the holder can exercise now;
  \item the European put, because the American holder has every European exercise choice plus the
        extra right to exercise early.
\end{itemize}

These two inequalities are not a complete proof of pricing accuracy, but they are essential
financial sanity checks. If either inequality fails materially, the American solver should not be
trusted for later boundary, Greek, stress-map, dataset, or neural-surrogate work.

\section{Validation Criteria}

\subsection{Obstacle Condition}

For the American put, the value grid should satisfy
\begin{equation}
  U_{n,i} \geq \Phi_i
\end{equation}
at every tested time and spot node. Ticket 05 reports
\begin{align}
  \min_{n,i}(U_{n,i}-\Phi_i), \\
  \max_{n,i}\max(\Phi_i-U_{n,i},0).
\end{align}
The first quantity is the minimum obstacle gap. The second is the maximum obstacle violation.

\subsection{American Versus European Put}

The American put should also satisfy
\begin{equation}
  U^{\mathrm{American}}(S,T)
  \geq
  P^{\mathrm{European}}(S,T)
\end{equation}
up to numerical tolerance. Ticket 05 compares final American values against the Ticket 01
European Black-Scholes put utility.

\subsection{Convergence Metadata}

Each PSOR time step returns metadata. Ticket 05 reports whether all steps converged, how many
PSOR iterations were needed, and the maximum final update. This matters because an obstacle check
alone can miss a solver that stopped too early.

\subsection{Selected Spots and Medium-to-Fine Comparison}

Selected spot checks use \(S/K\) values \(0.5\), \(0.8\), \(1.0\), \(1.2\), and \(1.5\). These
make the table readable for students because each row corresponds to a familiar moneyness point.
The medium-to-fine comparison is intentionally light: it compares selected prices from
\(M=N=120\) and \(M=N=180\). Full grid and domain sensitivity belongs to a later ticket.

\section{Obstacle Gap Is Not Boundary Extraction}

Ticket 05 reports the obstacle gap
\begin{equation}
  G(S,\tau)=U(S,\tau)-\Phi(S).
\end{equation}

This gap is useful because it directly checks the obstacle condition. If \(G<0\), the solver has
violated the American payoff constraint. If \(G=0\) at some grid points, those points may be
consistent with immediate exercise.

However, Ticket 05 does not extract a free boundary. Boundary extraction requires option-type
specific search direction, thresholds, interpolation, and missing-boundary metadata. That is
planned for a later ticket. In this report, obstacle gap means validation evidence only.

\section{Implementation Design Decisions}

\begin{longtable}{L{0.28\textwidth}L{0.60\textwidth}}
\toprule
Decision & Reason \\
\midrule
\endhead
Use Ticket 04 solver as-is &
Ticket 05 is a validation layer. Solver changes are out of scope unless a genuine bug is found. \\
Use \(K=1\), \(S_{\max}=4\), \(T=1\), \(r=0.05\), \(q=0.02\), \(\sigma=0.20\) &
These values match earlier ticket defaults and the validation plan's normalized-strike workflow. \\
Use medium and fine grids &
This gives a preliminary refinement check without turning Ticket 05 into full grid/domain
sensitivity. \\
Write CSV artifacts &
Tables are easy to inspect, test, and reuse in reports. \\
Create one figure &
An American put validation ticket benefits from one visual comparison of American value, payoff,
and European value. \\
Avoid boundary extraction language &
Obstacle gap reporting is allowed, but boundary location inference is explicitly out of scope. \\
\bottomrule
\end{longtable}

\section{Code-Level Function Explanations}

\begin{longtable}{P{0.38\textwidth}L{0.54\textwidth}}
\toprule
Function & Purpose \\
\midrule
\endhead
\code{run\_validation\_cases} &
Runs the medium and fine American put cases, returns validation rows, selected-spot rows, and
artifact metadata. \\
\code{summarize\_american\_put\_case} &
Builds one obstacle, European-comparison, and PSOR convergence summary row. \\
\code{selected\_spot\_rows} &
Reports payoff, European put, American put, American-minus-European, and American-minus-payoff at
selected moneyness points. \\
\code{nearest\_spot\_index} &
Selects the nearest grid point to a requested spot. It does not interpolate. \\
\code{write\_csv} &
Writes rows with stable column order. \\
\code{create\_value\_vs\_payoff\_figure} &
Creates the validation figure if matplotlib is available. \\
\code{main} &
Runs the experiment and writes the CSV/figure artifacts. \\
\bottomrule
\end{longtable}

\section{Validation Result Tables and Interpretation}

\subsection{Validation Cases}

{\small
\begin{longtable}{@{}L{0.16\textwidth}L{0.07\textwidth}L{0.07\textwidth}L{0.07\textwidth}L{0.07\textwidth}L{0.09\textwidth}L{0.07\textwidth}L{0.07\textwidth}@{}}
\toprule
Case & \(K\) & \(T\) & \(r\) & \(q\) & \(\sigma\) & \(M\) & \(N\) \\
\midrule
\endhead
Medium & 1.0 & 1.0 & 0.05 & 0.02 & 0.20 & 120 & 120 \\
Fine & 1.0 & 1.0 & 0.05 & 0.02 & 0.20 & 180 & 180 \\
\bottomrule
\end{longtable}
}

\subsection{Obstacle and European Comparison Summary}

{\small
\begin{longtable}{@{}L{0.13\textwidth}L{0.14\textwidth}L{0.16\textwidth}L{0.19\textwidth}L{0.19\textwidth}@{}}
\toprule
Case & Min gap & Max violation & Min Am. - Eur. & Max Am. - Eur. \\
\midrule
\endhead
Medium & 0 & 0 & \(-3.7241\times10^{-14}\) & 0.0487706 \\
Fine & 0 & 0 & \(-3.7241\times10^{-14}\) & 0.0487706 \\
\bottomrule
\end{longtable}
}

The minimum American-minus-European value is a tiny negative roundoff-scale value. It is not a
financially meaningful violation. The American put value dominates the European put across the
tested grid up to numerical tolerance.

\subsection{Convergence Metadata}

\begin{longtable}{@{}L{0.14\textwidth}L{0.18\textwidth}L{0.15\textwidth}L{0.18\textwidth}L{0.22\textwidth}@{}}
\toprule
Case & All steps converged & Step count & Max iterations & Max final update \\
\midrule
\endhead
Medium & True & 120 & 10 & \(9.9971\times10^{-9}\) \\
Fine & True & 180 & 10 & \(9.9919\times10^{-9}\) \\
\bottomrule
\end{longtable}

Every PSOR step converged in both validation cases. The maximum final updates are just below the
requested tolerance scale, which is expected for an iterative method that stops when the update
falls below tolerance.

\subsection{Selected Spot Price Comparison}

{\footnotesize
\begin{longtable}{@{}L{0.10\textwidth}L{0.10\textwidth}L{0.12\textwidth}L{0.13\textwidth}L{0.13\textwidth}L{0.13\textwidth}L{0.13\textwidth}@{}}
\toprule
Case & \(S/K\) & Payoff & European put & American put & Am. - Eur. & Am. - payoff \\
\midrule
\endhead
Medium & 0.5 & 0.500000 & 0.461146 & 0.500000 & 0.038854 & 0 \\
Medium & 0.8 & 0.200000 & 0.182378 & 0.200381 & 0.018003 & 0.000381 \\
Medium & 1.0 & 0 & 0.063301 & 0.066301 & 0.003001 & 0.066301 \\
Medium & 1.2 & 0 & 0.015602 & 0.015984 & 0.000381 & 0.015984 \\
Medium & 1.5 & 0 & 0.001228 & 0.001254 & 0.000026 & 0.001254 \\
Fine & 0.488889 & 0.511111 & 0.472032 & 0.511111 & 0.039079 & 0 \\
Fine & 0.8 & 0.200000 & 0.182378 & 0.200606 & 0.018228 & 0.000606 \\
Fine & 1.0 & 0 & 0.063301 & 0.066469 & 0.003168 & 0.066469 \\
Fine & 1.2 & 0 & 0.015602 & 0.016054 & 0.000451 & 0.016054 \\
Fine & 1.488889 & 0 & 0.001358 & 0.001386 & 0.000028 & 0.001386 \\
\bottomrule
\end{longtable}
}

The selected spot table shows the expected pattern: the American put is never below payoff and is
above the European put. At low spot values, the American value can equal payoff exactly on the
chosen grid node. This is an obstacle validation observation, not a boundary extraction claim.

\subsection{Medium-to-Fine Comparison}

The largest selected-spot medium-to-fine difference recorded by the experiment is \(0.0111111\).
This value is driven by nearest-grid selection at \(S/K=0.5\): the medium grid includes \(S=0.5\),
while the fine grid's nearest point is \(S=0.4888889\). This is a useful reminder that selected
spot tables using nearest nodes are simple and auditable, but they are not a full interpolation or
convergence study.

\section{Figure Interpretation}

Figure~\ref{fig:value-payoff} plots the fine-grid American put value, the payoff, and the
European put value over moneyness \(S/K\). The figure is evidence for three basic visual checks:
\begin{itemize}
  \item the American value does not fall below payoff;
  \item the American value is above the European put value;
  \item the curves have the expected put shape, with value decreasing as spot increases.
\end{itemize}

\begin{figure}[H]
\centering
\includegraphics[width=0.82\textwidth]{results/01_solver_validation/figures/ticket_05_american_put_value_vs_payoff.png}
\caption{Ticket 05 American put value, payoff, and European put comparison. The plot is a validation visual, not a boundary-extraction plot.}
\label{fig:value-payoff}
\end{figure}

\section{Testing Strategy}

\begin{longtable}{L{0.26\textwidth}L{0.32\textwidth}L{0.30\textwidth}}
\toprule
Test category & What it proves & What it does not prove \\
\midrule
\endhead
Script/artifact test &
The experiment runs and writes the expected CSV and figure files. &
Full financial validation across regimes. \\
CSV schema tests &
The result tables have stable, expected columns. &
That every metric is sufficient for later tickets. \\
Obstacle tests &
American values dominate payoff within tolerance. &
Complementarity residual accuracy. \\
European comparison tests &
American puts dominate European puts within tolerance. &
External benchmark agreement for American prices. \\
PSOR metadata tests &
Default cases converged and expose iteration/update summaries. &
Convergence for all possible grids and parameters. \\
Nearest-grid test &
Selected spot reporting is predictable. &
Interpolation accuracy. \\
Figure test &
The meaningful validation PNG is created when matplotlib is available. &
That the figure alone validates the solver. \\
Scope tests &
No boundary extraction, Greeks, stress maps, datasets, or neural APIs/files are introduced. &
Future tickets that intentionally add those features. \\
Docstring test &
New Python files identify Ticket 05. &
Completeness of the research report. \\
\bottomrule
\end{longtable}

\section{Test Results and Interpretation}

The test command was:
\begin{lstlisting}[language=bash]
PYTHONPYCACHEPREFIX=/private/tmp/codex_pycache PYTHONPATH=src MPLCONFIGDIR=/private/tmp/matplotlib-cache python3 -m unittest discover -s tests
\end{lstlisting}

Observed result during Ticket 05 implementation:
\begin{lstlisting}
Ran 50 tests

OK
\end{lstlisting}

The experiment command was:
\begin{lstlisting}[language=bash]
PYTHONPYCACHEPREFIX=/private/tmp/codex_pycache PYTHONPATH=src MPLCONFIGDIR=/private/tmp/matplotlib-cache python3 experiments/02_american_put_validation.py
\end{lstlisting}

The experiment wrote two validation rows, ten selected-spot rows, and one value-versus-payoff PNG.

\section{Implementation Pitfalls and Debugging Notes}

{\small
\begin{longtable}{@{}L{0.21\textwidth}L{0.24\textwidth}L{0.22\textwidth}L{0.19\textwidth}@{}}
\toprule
Recorded item or pitfall & Symptom & Fix or guard & Beginner lesson \\
\midrule
\endhead
Expected TDD red run &
The new tests failed with \code{FileNotFoundError} because the Ticket 05 experiment script did not
exist yet. &
Create \file{experiments/02_american_put_validation.py}. &
A red test proves the tests are aimed at the new artifact. \\
Relative metadata paths &
Two tests failed because the script returned relative artifact paths while tests expected absolute
paths. &
Root output paths at the project directory and return absolute metadata paths. &
Tests can reveal reproducibility details, not only finance logic. \\
Font cache warning &
Matplotlib emitted a fontconfig cache warning in the sandboxed runtime. &
Use writable \code{MPLCONFIGDIR} and \code{FC\_CACHEDIR} locations. &
Report-generation code may need environment hygiene. \\
Overclaiming from one case &
One representative case could be mistaken for full solver validation. &
State that this is preliminary American put validation only. &
Validation evidence has a scope. \\
Confusing obstacle gap with boundary extraction &
\(U-\Phi\) may tempt a reader to infer a free boundary. &
Report gap values but do not infer, label, or extract a boundary. &
Diagnostics and boundary extraction are different tasks. \\
Ignoring European comparison &
Checking only \(U\geq\Phi\) would miss a weak American value. &
Compare against the Ticket 01 European put utility. &
American value should include the European exercise right plus early exercise. \\
Ignoring PSOR non-convergence &
Obstacle values alone could hide slow or failed PSOR steps. &
Report all-step convergence, iterations, and final updates. &
Numerical credibility includes solver status metadata. \\
Using too coarse a grid &
A coarse grid can create misleading selected-spot values. &
Include medium and fine grids, and label this as a light check. &
Grid checks should be layered and honest. \\
\bottomrule
\end{longtable}
}

\section{Limitations}

\begin{longtable}{L{0.34\textwidth}L{0.54\textwidth}}
\toprule
Not included in Ticket 05 & Reason \\
\midrule
\endhead
Boundary extraction & Planned for a later ticket; this report does not infer a free-boundary
location. \\
Complementarity residual module & Planned diagnostic work after core validation layers. \\
Greeks & Delta and Gamma diagnostics are later work. \\
Stress maps, datasets, neural networks & These require a more complete validated solver pipeline. \\
No-dividend American call validation & Belongs to Ticket 06. \\
Dividend-paying American call validation & Belongs to a later call-validation ticket. \\
Full grid/domain sensitivity & The current medium-to-fine comparison is preliminary only. \\
External American reference prices & This ticket uses internal financial inequalities and solver metadata, not external benchmark
tables. \\
\bottomrule
\end{longtable}

\section{Lessons Learned / Study Notes}

The first lesson is that American put validation should not rely on one number. The value must be
checked against payoff, against the European put, and against PSOR convergence metadata.

The second lesson is that obstacle gap is a validation diagnostic. It helps confirm that the
American constraint is respected, but it should not be turned into a boundary claim without the
later extraction logic.

The third lesson is that selected-spot tables are excellent for learning because they translate a
surface into familiar moneyness points. They are not a replacement for full convergence and domain
sensitivity studies.

The fourth lesson is that figures should support an argument. The Ticket 05 figure shows the
relationship among American value, European value, and payoff. It is not decorative and does not
label a free boundary.

\section{Link to Ticket 06}

Ticket 06 should validate the no-dividend American call theorem. Under the standard
Black-Scholes assumptions with \(q=0\), a standard American call should match the corresponding
European call. That check is different from Ticket 05: it tests whether the American exercise
logic avoids inventing early exercise where theory says it should not occur.

\section{Reproducibility Record}

\subsection{Files Created}

\begin{longtable}{P{0.54\textwidth}L{0.38\textwidth}}
\toprule
Path & Role in Ticket 05 \\
\midrule
\endhead
\file{experiments/02_american_put_validation.py} &
Runs the American put validation cases and writes CSV/figure artifacts. \\
\file{tests/test_american_put_validation.py} &
Focused unittest coverage for Ticket 05 artifacts, metrics, scope, and docstrings. \\
\file{results/01_solver_validation/tables/ticket_05_american_put_validation.csv} &
Validation summary table. \\
\file{results/01_solver_validation/tables/ticket_05_american_put_selected_spots.csv} &
Selected moneyness price comparison table. \\
\file{results/01_solver_validation/figures/ticket_05_american_put_value_vs_payoff.png} &
American put value, payoff, and European put comparison figure. \\
\file{reports/03_solver/tickets/ticket_05_american_put_validation.tex} &
Formal LaTeX study-report source. \\
\file{reports/03_solver/tickets/ticket_05_american_put_validation.pdf} &
Formal PDF compiled from the LaTeX source. \\
\bottomrule
\end{longtable}

\subsection{Commands}

Test command:
\begin{lstlisting}[language=bash]
PYTHONPYCACHEPREFIX=/private/tmp/codex_pycache PYTHONPATH=src MPLCONFIGDIR=/private/tmp/matplotlib-cache python3 -m unittest discover -s tests
\end{lstlisting}

Experiment command:
\begin{lstlisting}[language=bash]
PYTHONPYCACHEPREFIX=/private/tmp/codex_pycache PYTHONPATH=src MPLCONFIGDIR=/private/tmp/matplotlib-cache python3 experiments/02_american_put_validation.py
\end{lstlisting}

Compile command:
\begin{lstlisting}[language=bash]
PYTHONPYCACHEPREFIX=/private/tmp/codex_pycache PYTHONPATH=src python3 -m compileall -q src tests experiments
\end{lstlisting}

LaTeX compile command:
\begin{lstlisting}[language=bash]
HOME=/private/tmp FC_CACHEDIR=/private/tmp/fontconfig-cache latexmk -xelatex -interaction=nonstopmode -halt-on-error -outdir=/private/tmp/ticket05_latex reports/03_solver/tickets/ticket_05_american_put_validation.tex
\end{lstlisting}

PDF verification uses \code{pdfinfo} and page rendering with \code{pdftoppm}. Raw Git status is
intentionally not included in this formal PDF report.

\bibliographystyle{plainnat}
\bibliography{reports/03_solver/references}

\end{document}
